% !TeX root = ../教材题库_全部条目_转写工作稿.tex
% Source: XJ-B1-C02-PR-P058-Y0432
% 第2章 / 2.3 一元二次不等式 / 教材 p.58 / PDF 第65页
% OCR未校对：公式与文字均可能有误；未完成数学TeX转写；禁止直接发给学生
% 原题及相邻原解截图（核对用）：../题目原图/XJ-B1-C02-PR-P058-Y0432-01.webp
% 本文件只在转写完成后作为题干使用；OCR 条目尚不是可靠题干。
\begin{bankitem}{p.58 练习 第2题}{第2章 / 2.3 一元二次不等式 / 教材 p.58 / PDF 第65页}{XJ-B1-C02-PR-P058-Y0432}{OCR未校对：公式与文字均可能有误；未完成数学TeX转写；禁止直接发给学生}
\begin{ocrdraft}
[教材p.58 / PDF第65页]
\par
2\ensuremath{\cdot}行驶中的汽车，由于具有惯性，刹车后还要继续向前滑行一段距离才能停
\par
住，这段距离称为“刹车距离”．刹车距离是分析交通事故的一个重要指标．
\par
在一个限速为40km/h的弯道上，甲、乙两辆汽车相向而行，发现情况不
\par
对，同时刹车，但还是相碰了．事后现场勘查测得甲车的刹车距离略超过12m，
\par
乙车的刹车距离略超过10m,又知甲、乙两种车型的刹车距离5（m)与车速
\par
禳km/h）分别满足如下关系式：
\par
=0．，十0．01t，2，、乙=0．05t，十0．00502．
\par
试用解一元二次不等式的方法判断甲、乙两辆汽车有无超速现象．
\end{ocrdraft}
\end{bankitem}
